\documentclass[10pt,a4paper]{amsart}
\usepackage[margin=19mm]{geometry}
\usepackage{amsmath,amssymb,mathtools}
\usepackage[scr=rsfs,cal=euler]{mathalfa}
\usepackage{xfrac}
\usepackage[unicode,hidelinks,bookmarks=false]{hyperref}
\newcommand{\ie}{i.e.\ }
\newcommand{\cf}{cf.\ }
\newcommand{\resp}{respectively\ }
\newcommand{\Subsection}{\subsection}
\providecommand{\nobreakdash}{\nobreak\hbox{-}}
\setlength{\emergencystretch}{2em}
\multlinegap0pt
\usepackage{amssymb}
\usepackage{mathtools}
\usepackage{xcolor}
\newtheorem{theorem}{Theorem}
\newcommand{\OmegaF}{\Omega_F}
\title{\textbf{A Nonmonotone Real-Rootedness Set for\\ Symmetric Imaginary Shifts}\\[0.7em] \large\textnormal{AI Research Artifact}}
\author{Vasily Stodolsky}
\date{Source: arXiv:2608.19240v1 (CC BY 4.0)}
\begin{document}
\maketitle

\noindent\emph{Source and licence.} \textbf{A Nonmonotone Real-Rootedness Set for\\ Symmetric Imaginary Shifts}\\[0.7em] \large\textnormal{AI Research Artifact}. Version arXiv:2608.19240v1 (CC BY 4.0), distributed by arXiv under the Creative Commons Attribution 4.0 International licence, http://creativecommons.org/licenses/by/4.0/, as recorded in the arXiv licence metadata for this exact version and shown on the abstract page https://arxiv.org/abs/2608.19240v1. Full licence notice: \texttt{licenses/arxiv-2608.19240-CC-BY-4.0.txt}.\par\smallskip

\noindent\textit{Editorial scope.} Counted unit: the article's theorem thm:main (source0.tex lines 176--192). The article's complete original proof of it is delivered with it (source0.tex lines 194--204, one \texttt{proof} environment of 11 lines). No mathematics is written, repaired or re-derived: the statement and the proof are the article's own lines, in the article's own notation, and no retained line is substituted or re-flowed.  Retained uncounted as the source-local context the proof needs: lines 158--167.  Nothing was omitted from any retained span, so the package carries no omission record.  No retained line contains a citation, so the wrapper carries no bibliography and no bibliography entry.  No retained line contains a figure environment, an image-inclusion command or drawing code, so no image is delivered and the package declares no asset dependency.  Editorial formatting only, in addition: an AMS \texttt{amsart} preamble that restates the article's own declarations and macros used by the retained lines, namely the environments \texttt{theorem} on separate counters, together with the standard packages those lines need.  The retained environments print in this document as Theorem 1 in order of appearance, whereas the article prints the same items with the numbering of its own full document.  The complete native source of the article as distributed in the arXiv e-print for this exact version is bundled unchanged as \texttt{sources/208072/source0.tex}.
\begin{align}
F(z)={}&
\left(\left(z-\frac1{50}\right)^2+\left(\frac7{25}\right)^2\right)
\left(\left(z+\frac1{50}\right)^2+\left(\frac7{25}\right)^2\right)
\notag\\
&\quad\times
\left(\left(z-\frac35\right)^2+\left(\frac{11}{25}\right)^2\right)
\left(\left(z+\frac35\right)^2+\left(\frac{11}{25}\right)^2\right).
\label{eq:F}
\end{align}

\begin{theorem}\label{thm:main}
The polynomial $F$ in \eqref{eq:F} is real and even, and every zero of $F$
lies in $|\operatorname{Im}z|\leq11/25$. Moreover,
\[
  \frac6{25}\in\OmegaF,
  \qquad
  \frac3{10}\notin\OmegaF,
  \qquad
  \frac{12}{25}\in\OmegaF.
\]
The respective numbers of distinct real zeros of $A_\omega$ are $8$, $4$,
and $8$. In addition,
\[
  \left[\frac{11}{25},\infty\right)\subset\OmegaF.
\]
Consequently $\OmegaF$ is neither an up-set nor an interval.
\end{theorem}

\begin{proof}[Proof of the symmetry and strip assertions]
The zeros of $F$ are visible directly from \eqref{eq:F}:
\[
 \frac1{50}\mathbin{\pm}\frac7{25}i,
 -\frac1{50}\mathbin{\pm}\frac7{25}i,
 \frac35\mathbin{\pm}\frac{11}{25}i,
 -\frac35\mathbin{\pm}\frac{11}{25}i.
\]
This proves the symmetry and strip assertions. The two remaining parts of
the theorem are proved separately below.
\end{proof}

\end{document}
